\section{Introduction}

The \BA\ (BA) function is a central object in the algebro-geometric
integration of nonlinear equations. Krichever's construction, building
on earlier work of Baker and Akhiezer, associates such functions with
an algebraic curve, marked points and prescribed divisor data
\cite{Krichever1977,KricheverShiota2011}. They solve the auxiliary
linear problems underlying the Kadomtsev--Petviashvili hierarchy and,
with the appropriate spectral data, finite-gap Korteweg--de Vries
solutions. The same functions enter the integration of
finite-dimensional systems, including Neumann and Garnier systems;
Mumford's polynomial coordinates describe the associated moving
divisors \cite{Kostov1999,Eilbeck2000,Mumford1984,Vanhaecke2001}.

These constructions connect functions and dynamics in two directions.
A Lax representation associates spectral curves and eigenvector data
with a Hamiltonian system. Conversely, prescribed curve and divisor
data produce BA functions from which solutions can be constructed.
Adler--Kostant--Symes theory explains the commuting Hamiltonians in
many matrix realizations, while spectral line bundles describe their
algebraic integration \cite{AdlerVanMoerbeke2018,Beauville1990}.
Together these methods provide a broad framework, within which
particular systems select particular curves, marked points and
functional realizations.

Finite-dimensional multiwave equations provide a direct setting for
this relationship between functions and dynamics. In genus one,
elliptic sigma ratios describe individual modes of integrable
three-wave and four-wave interactions
\cite{Hesketh2015,HeskethFWM2026}; their shifts and exponential factors
encode both powers and phases. Multiwave systems also have a history
of Hamiltonian and spectral constructions that extends beyond these
elliptic examples.

Three-wave interaction equations arise in models of ocean waves,
plasma waves and quadratic nonlinear optical media. Restricted
multiple three-wave systems couple several triads through a common
mode \cite{KostovTsiganov2008}; continuous-wave optical interactions
also lead to finite-dimensional mode equations \cite{Hesketh2015}.
Menyuk, Chen and Lee investigated integrable couplings through
Hamiltonian, Lax and Painlev\'e analyses
\cite{MenyukChenLee1982,MenyukChenLee1983Integrable,
MenyukChenLee1983Painleve}. They obtained a complete commuting
integral set for the equal-coupling shared-mode family, constructed
broader wave networks, and integrated the equal-mismatch family
using elliptic functions and phase quadratures.
Rauch-Wojciechowski, Jiang and Bullough gave an alternative matrix
Lax representation and constructed a broader class of multiwave
ordinary differential equations \cite{WJB1986}. Jiang and
Rauch-Wojciechowski then developed this approach in a graded Lie
algebra framework, constructing canonical cubic systems with
rectangular mode arrays and studying the completeness of their
commuting integrals \cite{JRW1991}.
Kostov and Tsiganov studied the restricted shared-mode family through
a two-component Dirac problem. They constructed its hyperelliptic
spectral curve and finite-gap eigenfunction representation, together
with a classical $r$-matrix and bi-Hamiltonian structure; the
single-triad case was evaluated explicitly in Weierstrass functions,
including its phase \cite{KostovTsiganov2008}.

\subsection{Our contribution}

Our starting point is the function itself. Given a smooth curve with
a marked point and finitely many scalar one-point BA evaluations, we
ask how to place them in a canonical polynomial Hamiltonian system.
Which additional modes are needed, and how do the curve-point data
determine the coupling weights and linear phase rates? We answer
this realization question within the canonical multiwave family
developed in these earlier works.

The vector--matrix--vector member provides a canonical Hamiltonian
realization of the scalar BA function and its Abelian evolution.
The correspondence works in both directions: generic initial modes
reconstruct BA evaluations and their pair functions, while any finite
collection of regular evaluations on any smooth curve with a marked
point can be included after adjoining points. The prescribed curve
is recovered as the normalization of the resulting spectral level.
Moreover, the commuting spectral Hamiltonians generate all Abel
directions, with the accompanying mode phases determined explicitly.
The curve type changes the compatible phase coefficients, initial
values and required mode counts, but the cubic Hamiltonian retains
its form. We call this $N=1$ realization the \emph{canonical multiwave
Baker--Akhiezer system} (CMBAS). Its two-way correspondence and
Hamiltonian hierarchy make the system a common canonical setting
for scalar BA dynamics across curve types.

We write the multiwave modes as three rectangular matrices
$\boldsymbol u\in\C^{L\times M}$,
$\boldsymbol v\in\C^{M\times N}$ and
$\boldsymbol w\in\C^{N\times L}$, where $L,M,N$ are positive
integers. Their entries are indexed by
$\ell=1,\ldots,L$, $m=1,\ldots,M$ and $n=1,\ldots,N$.
Each entry has an independent Hamiltonian partner, denoted by a bar.
Stars denote actual complex conjugation, $\trans$ denotes transpose,
and $\dagger$ denotes actual conjugate transpose. The barred matrices
have the same dimensions as their unbarred partners. On the all-positive
real form $\bar{\boldsymbol u}=\boldsymbol u^*$, and similarly for
$\boldsymbol v,\boldsymbol w$.
Thus the cubic mixing terms have the form
\begin{align}\label{eq:intro-mixing}
 H_{\mathrm{mix}}
 &=\kappa\sum_{\ell=1}^{L}\sum_{m=1}^{M}\sum_{n=1}^{N}
 \left(u_{\ell m}v_{mn}w_{n\ell}
       +\bar{u}_{\ell m}\bar{v}_{mn}\bar{w}_{n\ell}\right)\nonumber\\
 &=\kappa\,\tr\!\left(
 \boldsymbol u\boldsymbol v\boldsymbol w
 +\bar{\boldsymbol u}\bar{\boldsymbol v}\bar{\boldsymbol w}
 \right),
\end{align}
with constant coupling $\kappa$. Each interaction joins one mode from
each matrix; quadratic terms supply the linear phase rates. For
$N=1$, the $v$ and $w$ matrices become two vectors. In our realization,
ordinary and dual scalar BA evaluations supply these elementary
vectors, while a normalized two-point BA function supplies the
pair-mode matrix $u$.
The pair modes are independent canonical variables of the Hamiltonian
system. The full three-matrix family supplies the parent Hamiltonian and
spectral framework. We establish the complete Liouville count for
all sizes and phase profiles in \cref{thm:canonical-independence},
with its independence proof in Appendix~\ref{sec:integrability-proof}.
The main BA results concern $N=1$ and are as follows.

\begin{itemize}
\item \textbf{A necessary and sufficient BA realization criterion.}
For $N=1$, \cref{thm:general-BA-realization} determines precisely when
ordinary/dual BA evaluations and their pair kernels satisfy the
canonical equations on an open Jacobian family. Meromorphic
interpolation determines the weights and phases through $g$ linear
residue equations on a genus-$g$ curve. Generic groups of at least
$g$ points admit the required weights; special smaller groups obey
the same consistency test. Every finite list of regular elementary
evaluations can be included after adjoining points when necessary.
A quadratic kernel residue identity supplies the Hamilton equations.
The explicit spectral-coordinate
map in \cref{prop:BA-prescribed-spectral-curve} proves that the
prescribed curve normalizes the realized Lax spectral curve.

\item \textbf{BA reconstruction from generic canonical initial modes.}
Conversely, \cref{thm:BA-spectral-converse} reconstructs scalar
one-point BA solutions from generic canonical initial data of the
$N=1$ system. The spectral curve, eigenline divisor and marked points
recover the elementary and pair modes. Repeated phases are treated
by constant canonical block bases and active reduction; the original
modes can then be linear combinations of the reconstructed BA
evaluations.

\item \textbf{Higher BA times and constant mode phases.}
\Cref{thm:BA-higher-flows} identifies the commuting spectral flows on
every prescribed BA realization. A meromorphic normalization computes
their one-point BA times and constant reciprocal mode phases; the full
spectral family spans every Jacobian direction on each fixed spectral
level. The required linear combinations depend on the curve data. The construction also
applies when phases coincide or the plane spectral image is singular.
\Cref{cor:BA-higher-reconstruction} extends the generic reconstruction
to these higher flows using the same curve and reference BA data.

\end{itemize}

The hyperelliptic specialization makes the realization criterion
explicit through polynomial interpolation in Baker sigma and
Kleinian conventions. It includes genus-one realizations with
arbitrarily many elementary elliptic modes and their cubic pair
interactions. The mode-power map to Mumford coordinates connects
the construction to established divisor dynamics while retaining
the individual BA modes for phase reconstruction. Numerical examples
on a genus-two hyperelliptic curve and a genus-three trigonal curve
compare the sigma formulas with independent integration of the
canonical equations.

The independent phase parameters used here already fit the earlier
multiwave framework. Our realization determines their compatible
values from the BA curve-point data. Appendices~\ref{sec:MCL-dictionary}
and~\ref{sec:JRW-dictionary} give the exact comparisons with
Menyuk--Chen--Lee and Jiang--Rauch-Wojciechowski.
The marked spectral-curve description explains the different covering
types and the distinction between regular and special levels.
The paper then develops the BA correspondence and its commuting flows,
followed by the explicit hyperelliptic interpolation, Mumford projection
and numerical validation.
